% EGA V, English translation by P. Blass, 1987 (cote 24): archive cote 24, site page 51 (PDF page 52).
% typescript. Draft compared: none.
% AI-transcribed by Claude Opus 5.5 (claude-opus-5-5, adaptive thinking, effort high), Message Batches API, batch ega24b (msgbatch_01RGJJeHdR96U16ZWFuUjtDx), 2026-11-13; not yet checked by eye.
\scanpage{24/51}
\mspage{51}

thus that the inverse image of $\underline{Y}^{1}$ (which is open in the above
closure) by h is an open subset of U, let us call it U1.  Because
of hg$_{1}$ = inclusion we see immediately that g$^{-1}$(U) $\to$ U is an open
immersion.  Since X is proper over k i.e. closed in $\underline{P}$, therefore g
is proper therefore surjective thus g$_{1}$ : g$^{-1}$(U) $\to$ U is
surjective, thus an isomorphism.  Also if W is an open subset of T
such that g$^{-1}$(W) $\to$ W is an isomorphism it follows that W is smooth
since $\underline{Y}^{1}$ is smooth, thus W \hand{$\subset$} U.  This proves 8.9.

The final assertions of 8.8 $\underline{Y}$sup sing = $\emptyset$ or dim $\underline{Y}$sup sing =
r - 2 and dim T = r - 1 are trivial:  the first one follows the
fact that $\underline{Y}$sing is irreducible of dim r and from the fact that Ysing
or Ysup sing [illegible]* is defined by the vanishing of a section D
of an invertible module.; the second from the fact that L being
finite over K we have deg tr L/k = deg tr K/k i.e. dim T = dim
$\underline{Y}$sing = r - 1.

\uline{Remark 8.10}

As we remarked in 8.9 with the notations of the corollary we
have g$^{-1}$(U) \hand{$\subset$} $\underline{Y}$sing - $\underline{Y}$sup sing but we notice that even if X is
closed in $\underline{P}$ this inclusion is not necessarily an equality, in
other words (noting that g$^{-1}$(U) is nothing else but the set of
points where g is étale, so that $\underline{Y}$sing - $\underline{Y}$sup sing is the set of
points where g is unramified) it may happen that there are points
y $\in$ $\underline{Y}$sing where g is unramified but not étale (which implies in
addition that g(y) is a point that is not geometrically normal and
even not geometrically unibranch of T).  In geometric terms this
corresponds to the following phenomenon; we may have a tangent

\margin{* ASK AG}
